\documentclass[10pt,a4paper]{amsart}
\usepackage{amsmath,amssymb,amsthm,graphicx}
% Retained preamble for the Stacks Project chapter snapshot in this directory.
% The chapter's own preamble.tex is a copy of the shared Stacks Project preamble; the block
% below is sliced MECHANICALLY, line for line, from it. Nothing was added, removed or
% redefined; the wrapper class, the xr-hyper/hyperref wiring and the externaldocument list
% are omitted because the assembler supplies the document class.
%
% lines taken from preamble.tex: 3, 5, 9, 11, 12, 156, 157, 158, ... (51 lines total)
%
% package coverage restored from the chapter's own preamble (2026-10-13): xy, xy-2cell

\usepackage[all]{xy}
\xyoption{2cell}
\UseAllTwocells

\newenvironment{reference}{\comment}{\endcomment}
\newenvironment{slogan}{\comment}{\endcomment}
\newenvironment{history}{\comment}{\endcomment}
\theoremstyle{plain}
\newtheorem{theorem}[subsection]{Theorem}
\newtheorem{proposition}[subsection]{Proposition}
\newtheorem{lemma}[subsection]{Lemma}
\theoremstyle{definition}
\newtheorem{definition}[subsection]{Definition}
\newtheorem{example}[subsection]{Example}
\newtheorem{exercise}[subsection]{Exercise}
\newtheorem{situation}[subsection]{Situation}
\theoremstyle{remark}
\newtheorem{remark}[subsection]{Remark}
\newtheorem{remarks}[subsection]{Remarks}
\numberwithin{equation}{subsection}
\def\lim{\mathop{\mathrm{lim}}\nolimits}
\def\colim{\mathop{\mathrm{colim}}\nolimits}
\def\Spec{\mathop{\mathrm{Spec}}}
\def\Hom{\mathop{\mathrm{Hom}}\nolimits}
\def\Ext{\mathop{\mathrm{Ext}}\nolimits}
\def\SheafHom{\mathop{\mathcal{H}\!\mathit{om}}\nolimits}
\def\SheafExt{\mathop{\mathcal{E}\!\mathit{xt}}\nolimits}
\def\Sch{\mathit{Sch}}
\def\Mor{\mathop{\mathrm{Mor}}\nolimits}
\def\Ob{\mathop{\mathrm{Ob}}\nolimits}
\def\Sh{\mathop{\mathit{Sh}}\nolimits}
\def\NL{\mathop{N\!L}\nolimits}
\def\CH{\mathop{\mathrm{CH}}\nolimits}
\def\proetale{{pro\text{-}\acute{e}tale}}
\def\etale{{\acute{e}tale}}
\def\QCoh{\mathit{QCoh}}
\def\Ker{\mathop{\mathrm{Ker}}}
\def\Im{\mathop{\mathrm{Im}}}
\def\Coker{\mathop{\mathrm{Coker}}}
\def\Coim{\mathop{\mathrm{Coim}}}
\DeclareMathSymbol{\boxtimes}{\mathbin}{AMSa}{"02}
\def\QCohstack{\mathcal{QC}\!\mathit{oh}}
\def\Cohstack{\mathcal{C}\!\mathit{oh}}
\def\Spacesstack{\mathcal{S}\!\mathit{paces}}
\def\Quotfunctor{\mathrm{Quot}}
\def\Hilbfunctor{\mathrm{Hilb}}
\def\Curvesstack{\mathcal{C}\!\mathit{urves}}
\def\Polarizedstack{\mathcal{P}\!\mathit{olarized}}
\def\Complexesstack{\mathcal{C}\!\mathit{omplexes}}
\def\Pic{\mathop{\mathrm{Pic}}\nolimits}
\def\Picardstack{\mathcal{P}\!\mathit{ic}}
\def\Picardfunctor{\mathrm{Pic}}
\def\Deformationcategory{\mathcal{D}\!\mathit{ef}}
\usepackage{cleveref}
\crefname{originallemma}{original lemma}{original lemmas}
\Crefname{originallemma}{Original lemma}{Original lemmas}
\makeatletter
\newlabel{lemma-characterize-versal}{{07YM}{source}{}{sourcepointer-1}{}}
\newlabel{lemma-characterize-versal@cref}{{[originallemma][0][]07YM}{[1][0][]source}{}{}{}}
\makeatother
\crefname{originallemma}{original lemma}{original lemmas}
\Crefname{originallemma}{Original lemma}{Original lemmas}
\makeatletter
\newlabel{more-algebra-lemma-cut-complex-in-two}{{068U}{source}{}{sourcepointer-2}{}}
\newlabel{more-algebra-lemma-cut-complex-in-two@cref}{{[originallemma][0][]068U}{[1][0][]source}{}{}{}}
\makeatother
\crefname{originallemma}{original lemma}{original lemmas}
\Crefname{originallemma}{Original lemma}{Original lemmas}
\makeatletter
\newlabel{remark-canonical-element}{{07YC}{source}{}{sourcepointer-3}{}}
\newlabel{remark-canonical-element@cref}{{[originallemma][0][]07YC}{[1][0][]source}{}{}{}}
\makeatother
\title{Openness of versality from a perfect obstruction theory input}
\author{The Stacks Project Authors}
\date{stacks-project:artin.tex:local-snapshot-20261012; GNU FDL 1.2 or later}
\begin{document}
\maketitle
\noindent\textit{Editorial note.} Statement span is lines 4363-4385 of sources/stacks-artin/pinned-original-source.tex; proof span is lines 4387-4402, the complete author proof (14 body lines).\par\medskip
\noindent\textit{Editorial note.} Navigation stubs: 3 cited-but-not-included labels are supplied as non-load-bearing pointers with their pinned tags where the snapshot tag index has them.\par\medskip
\noindent\textit{Editorial note.} extra\_packages is empty on purpose: the retained preamble already loads xy with options and a second declaration clashes in the assembler wrapper.\par\medskip
\noindent\textit{Editorial note.} Text is retained verbatim, including the author's own typos and placeholders.\par\medskip
\noindent\textit{Editorial note.} Queue index 42 of queue\_artin.json (pb=631, stmt 4363-4385, proof 4387-4402).\par\medskip
\noindent\textit{Editorial note.} stacks-artin wave ar01 (2026-10-14): fresh adjudication accept per lane-s3/artin\_judgments.md; pool per Chief 10:5x direct dispatch (idle-first over S12's 01:5x note, SMC precedent); artin.tex 99 rows/67 with proof judged 46 accept / 21 reject / 5 AGP-overlap avoidance rows; patch-sd-envs contract spans scanned.\par\medskip
\noindent\textit{Allowed standard prerequisites (editorial).}\begin{itemize}
\item Explicit prerequisite boundary. The cited inputs are prerequisites; the argument in the retained span is not replaced by any cited result.
\end{itemize}\medskip
\ifdefined\hypertarget\hypertarget{sourcepointer-1}{}\fi
\noindent\textit{Source reference.} Original lemma 07YM: sources/stacks-artin/pinned-original-source.tex, source lines 4235--4251. Cited by the retained statement or proof but not part of the collected spans; supplied only as a navigation pointer. The argument does not depend on its content.\par\medskip
\ifdefined\hypertarget\hypertarget{sourcepointer-2}{}\fi
\noindent\textit{Source reference.} Original lemma 068U: sources/stacks-artin/pinned-original-source.tex, source lines 4402--4402. Cited by the retained statement or proof but not part of the collected spans; supplied only as a navigation pointer. The argument does not depend on its content.\par\medskip
\ifdefined\hypertarget\hypertarget{sourcepointer-3}{}\fi
\noindent\textit{Source reference.} Original lemma 07YC: sources/stacks-artin/pinned-original-source.tex, source lines 3806--3838. Cited by the retained statement or proof but not part of the collected spans; supplied only as a navigation pointer. The argument does not depend on its content.\par\medskip
\begin{lemma}
\label{lemma-openness}
Let $S$ be a locally Noetherian scheme. Let $\mathcal{X}$ be a category
fibred in groupoids over $(\Sch/S)_{fppf}$ satisfying (RS*).
Let $U = \Spec(A)$ be an affine scheme of finite type over $S$ which maps
into an affine open $\Spec(\Lambda)$. Let $x$ be an object of $\mathcal{X}$
over $U$. Let $\xi : E \to \NL_{A/\Lambda}$ be a morphism of $D^{-}(A)$.
Assume
\begin{enumerate}
\item[(i)] for every deformation situation $(x, A' \to A)$ we have:
$x$ lifts to $\Spec(A')$ if and only if
$E \to \NL_{A/\Lambda} \to \NL_{A/A'}$ is zero,
\item[(ii)] there is an isomorphism of functors
$T_x(-) \to \Ext^0_A(E, -)$
such that $E \to \NL_{A/\Lambda} \to \Omega^1_{A/\Lambda}$
corresponds to the canonical element (see
Remark \ref{remark-canonical-element}),
\item[(iii)] the cohomology groups of $E$ are finite $A$-modules.
\end{enumerate}
If $x$ is versal at a closed point $u_0 \in U$,
then there exists an open neighbourhood $u_0 \in U' \subset U$
such that $x$ is versal at every finite type point of $U'$.
\end{lemma}

\begin{proof}
Let $C$ be the cone of $\xi$ so that we have a distinguished triangle
$$
E \to \NL_{A/\Lambda} \to C \to E[1]
$$
in $D^{-}(A)$. By Lemma \ref{lemma-characterize-versal}
the assumption that $x$ is versal at $u_0$ implies that
$H^{-1}(C \otimes^\mathbf{L} k) = 0$. By
More on Algebra, Lemma \ref{more-algebra-lemma-cut-complex-in-two}
there exists an $f \in A$ not contained in the prime corresponding to
$u_0$ such that $H^{-1}(C \otimes^\mathbf{L}_A M) = 0$ for
any $A_f$-module $M$. Using
Lemma \ref{lemma-characterize-versal}
again we see that we have versality for all finite type points of
the open $D(f) \subset U$.
\end{proof}

\end{document}
